\chapter{Recadrage et effets géométriques}
L’onglet \emph{Géométrie}, en haut du panneau de droite, regroupe l’orientation,
le recadrage et les transformations manuelles. Rien n’est écrit dans la photo :
la géométrie s’enregistre dans son compagnon et s’annule comme les autres étapes.
Tant que l’onglet est ouvert, l’aperçu montre tout le cadre, le recadrage sous
forme de rectangle ; revenir à l’onglet \emph{Réglages} affiche l’image recadrée.

\section{Orientation}
\emph{Pivoter de 90° à droite}, \emph{Miroir horizontal} et \emph{Miroir vertical}
ne dégradent aucun pixel. Le recadrage et les transformations suivent l’image.

\section{Redresser d’après les lignes}
IPR lit les lignes droites de la photo et propose une correction, que les curseurs
de \emph{Transformations} montrent ensuite et qui s’annule comme les autres étapes.
\emph{Niveau} propose une rotation ; \emph{Verticales} rend aussi les verticales
parallèles, si la photo en a assez ; \emph{Complet} fait de même avec les horizontales.
La proposition vaut surtout pour un horizon ou un bâtiment. Sans lignes assez longues
ou qui s’accordent, IPR n’en fait aucune : le bouton se borde de rouge et grossit
légèrement, et un message \emph{Erreur} donne la raison pendant deux secondes avant de
s’estomper. Tout bouton dont la commande échoue réagit de la même façon, et la raison
reste écrite dans la zone des messages, en bas de la fenêtre.

\attention{IPR ne connaît pas la scène : deux fils électriques en travers d’un paysage
lui ressemblent à un horizon. Regarder le résultat avant de le garder.}

\section{Recadrage}
Déplacer le rectangle ou l’une de ses huit poignées. L’image n’est recalculée
qu’au relâchement, et le geste compte pour une étape. La liste \emph{Format} impose
des proportions, y compris celles d’origine ; les champs \emph{Gauche}, \emph{Haut},
\emph{Largeur} et \emph{Hauteur} donnent une saisie numérique en pour cent du cadre.

\section{Transformations}
\begin{itemize}
\item \emph{Rotation} : de -45° à 45°, pour redresser un horizon. Une grille
d’alignement apparaît dans le cadre tant qu’un curseur de transformation est tenu ;
\emph{Toujours afficher la grille} la garde visible.
\item \emph{Vertical} et \emph{Horizontal} : corrigent les lignes fuyantes. Une valeur
positive élargit le haut, ou la droite.
\item \emph{Aspect}, \emph{Échelle}, \emph{Décalage horizontal} et \emph{Décalage
vertical} : ajustent les proportions, la taille et la place de l’image dans le cadre.
\end{itemize}
Ces réglages forment une seule transformation, appliquée en une seule passe : les
cumuler ne dégrade pas davantage l’image. Elle ne corrige pas la distorsion d’un
objectif, qui reste le rôle des profils de correction.

\section{Zones vides}
Une transformation peut laisser une partie du cadre sans image. IPR n’invente rien
pour la remplir : ces zones restent transparentes en PNG et TIFF, et noires en JPEG
et DNG, ce que le message d’export signale. Pour les retirer, recadrer, ou cocher
\emph{Contraindre à l’image}, qui réduit le rectangle autour de son centre jusqu’à
ce qu’il ne contienne que l’image.

\attention{Copier les paramètres, ou les appliquer à la sélection, ne transmet pas
la géométrie : le cadrage appartient à chaque photo. \emph{Réinitialiser les
réglages} la retire en revanche, comme les autres retouches.}
